% ECONOMICS_PROBLEM: {"id": "C65-29", "title": "Nonconvex quadratic BSDE uniqueness: a necessary correction", "jel_code": "C65", "source_id": "OP-0566"}
% FIRST_POSTED: 2026-10-10
% ATTEMPT_STATUS: solved
\documentclass[11pt]{article}
\usepackage[a4paper,margin=25mm]{geometry}
\usepackage{amsmath,amssymb,amsthm,microtype,titling}
\usepackage[hidelinks]{hyperref}
\setlength{\droptitle}{-12mm}
\newtheorem{theorem}{Theorem}
\newcommand{\E}{\mathbb E}
\newcommand{\Pp}{\mathbb P}
\title{Nonuniqueness of a quadratic BSDE\protect\\ with all exponential moments}
\author{Ji Ho Bae}
\date{10 October 2026}
\hypersetup{pdftitle={Nonuniqueness of a quadratic BSDE with all exponential moments},pdfauthor={Ji Ho Bae}}
\begin{document}
\maketitle
\begin{abstract}
We construct a scalar quadratic BSDE with a predictable generator, independent of the value variable and satisfying a uniform quadratic Lipschitz bound, that has a continuum of distinct solutions with the same terminal value. Each constructed solution has a square-integrable integrand and an integrable exponential of its running absolute maximum at every positive order. The generator is neither convex nor concave. Thus the nonconvex uniqueness question remains false even after imposing Lipschitz continuity in the value variable.
\end{abstract}

Problem~5.5 of Fan, Hu and Tang~\cite{FHT} asks about uniqueness for nonconvex quadratic BSDEs with all exponential moments. Its literal formulation does not control the dependence on $y$. The construction below is independent of $y$, so it also answers the repaired question that imposes a uniform Lipschitz bound in $y$ and a bounded generator at the origin. Class (D) means uniform integrability over all stopping times bounded by $T$.

\begin{theorem}\label{main}
For every $T>0$, on the usual augmented natural filtration of a one-dimensional Brownian motion $B$, there exist a predictable generator $g(t,\omega,z)$ and a terminal value $\xi$ such that, for all $z,z'\in\mathbb R$,
\begin{equation}\label{bounds}
 g(t,0)=0,\qquad |g(t,z)|\leq z^2,\qquad
 |g(t,z)-g(t,z')|\leq4(1+|z|+|z'|)|z-z'|.
\end{equation}
Every function $g(t,\omega,\cdot)$ is continuous, neither convex nor concave. The equation
\begin{equation}\label{bsde}
 Y_t=\xi+\int_t^T g(s,Z_s)\,ds-\int_t^T Z_s\,dB_s
\end{equation}
has a solution $(Y^c,Z^c)$ for every $c\in[0,1]$, with distinct initial values, and
\begin{equation}\label{moments}
 \mathbb E\int_0^T|Z^c_s|^2ds<\infty,\qquad
 \mathbb E\exp\!\left(\mu\sup_{0\leq t\leq T}|Y^c_t|\right)<\infty
 \quad\text{for every }\mu>0.
\end{equation}
In particular, $\mathbb E e^{\mu|\xi|}<\infty$ and $e^{\mu|Y^c|}$ is of class (D) for every $\mu>0$.
\end{theorem}

\paragraph{A sequence of stopped excursions.}
Put $s_k=T(1-2^{-k})$ for $k\geq0$. Leg $k\geq1$ occupies $[s_{k-1},s_k)$, has sign $\sigma_k=(-1)^{k+1}$, and belongs to cycle $n=\lceil k/2\rceil$. Set
\[
 a_n=\log(n+1),\qquad h_k(t)=\frac1{s_k-t},\qquad
 L^k_t=\sigma_k\int_{s_{k-1}}^t h_k(s)\,dB_s
            -\int_{s_{k-1}}^t h_k(s)^2ds.
\]
Extend $h_k$ by zero outside its leg. Let $\tau_k$ be the first exit of $L^k$ from $(-a_n,a_n)$. The deterministic clock $\int_{s_{k-1}}^t h_k^2$ tends to infinity as $t\uparrow s_k$. In clock time $L^k$ is a Brownian motion with drift $-1$, so $\tau_k<s_k$ almost surely. Call the leg successful if $L^k_{\tau_k}=a_n$.

Only legs whose predecessors were all successful are used. Write $A_k$ for this event and define the predictable process
\begin{equation}\label{q}
 q_t=\sum_{k\geq1}\mathbf1_{A_k}\mathbf1_{\{s_{k-1}<t\leq\tau_k\}}h_k(t),
 \qquad
 X_t=\int_0^t q_s\,dB_s-\int_0^t\sigma(s)q_s^2ds,
\end{equation}
where $\sigma(s)=\sigma_k$ on leg $k$ and endpoint values are immaterial. The event $A_k$ is known at time $s_{k-1}$, and the stopped interval indicators are predictable. Finiteness of the energy in~\eqref{q} will be verified below; equivalently, one may first construct these processes up to any finite number of legs.

The first leg in a cycle starts at $X=0$ and, on success, ends at $X=a_n$. The second starts at $a_n$ and, on success, returns to zero. A failed first leg ends at $-a_n$; a failed second leg ends at $2a_n$. After a failure, $q=0$ forever. Between a successful exit and the next leg, $X$ is constant.

For Brownian motion with drift $-1$ started at zero in $(-a,a)$, the probability of exiting at $a$ is
\[
 p(a)=\frac1{1+e^{2a}}.
\]
Indeed, apply It\^o's formula and bounded optional sampling to the harmonic function
$\frac{e^{2(x+a)}-1}{e^{4a}-1}$. The exit time $S$ in clock time also satisfies $\mathbb E S\leq a$: at $S\wedge m$, the stopped process lies in $[-a,a]$ and has expectation $-\mathbb E(S\wedge m)$, and then monotone convergence applies.

Thus each used leg in cycle $n$ has conditional success probability
$p_n=(1+(n+1)^2)^{-1}$. Independence of Brownian increments on the successive deterministic intervals gives, for the cycle $N$ of the first failure,
\begin{equation}\label{tails}
 \mathbb P(N\geq n)=\prod_{j<n}p_j^2\leq\frac1{(n!)^4},\qquad
 \mathbb E\int_0^Tq_t^2dt\leq2\sum_{n\geq1}\frac{a_n}{(n!)^4}<\infty.
\end{equation}
In particular $N<\infty$ almost surely, and
\begin{equation}\label{maximum}
 \sup_{t\leq T}|X_t|\leq2a_N,\qquad
 \mathbb E e^{\mu\sup|X|}\leq\sum_{n\geq1}\frac{(n+1)^{2\mu}}{(n!)^4}<\infty.
\end{equation}
All paths under consideration stop strictly before $T$. The null event of infinitely many successes can be discarded using the completed filtration.

\paragraph{A bounded perturbation with zero terminal value.}
Define
\[
 F_n(x)=\frac{1-e^{-2(x+a_n)}}{1-e^{-4a_n}},\qquad
 r_n=F_n(0)=\frac1{1+e^{-2a_n}},\qquad
 d_k=\prod_{j=k}^{\infty}r_{\lceil j/2\rceil},\qquad \eta=\frac1{16}.
\]
These products are positive: $\sum_n\log(1+(n+1)^{-2})<\infty$. They satisfy $d_k=r_n d_{k+1}$ when $n=\lceil k/2\rceil$. On $[-a_n,a_n]$,
\begin{equation}\label{F}
 F_n(-a_n)=0,\quad F_n(a_n)=1,\quad
 F_n''=-2F_n',\quad 0<F_n'\leq\frac{32}{15}.
\end{equation}
During a used leg $k$ of cycle $n$, up to its exit, set
\[
 D_t=\eta d_{k+1}F_n(L^k_t),\qquad
 K_t=\eta d_{k+1}\sigma_k h_k(t)F_n'(L^k_t).
\]
Hold $D$ constant between legs, set it to zero forever after failure, and put $K=0$ outside the used portions of legs. At the beginning of leg $k$, $D=\eta d_k$; on success it becomes $\eta d_{k+1}$. Thus the pieces join continuously. We have
\begin{equation}\label{D}
 D_0=\eta d_1>0,\qquad D_T=0,\qquad
 0\leq D\leq\eta,\qquad |K_t|\leq\frac{2}{15}q_t.
\end{equation}
It\^o's formula, using the drift $-1$ of $L^k$ and~\eqref{F}, gives the global identity
\begin{equation}\label{linear}
 dD_t=K_t\,dB_t-2\sigma(t)q_tK_t\,dt.
\end{equation}
All terms are integrable by~\eqref{tails} and~\eqref{D}. This construction uses no change of probability measure.

\paragraph{The generator and the solutions.}
Let
\[
 G(u)=2\bigl(1-|u-\tfrac32|\bigr)_+.
\]
This tent function is supported on $[1/2,5/2]$, is $2$-Lipschitz, satisfies $0\leq G(u)\leq u^2$, and equals $2u-1$ on $[1/2,3/2]$. Define
\begin{equation}\label{generator}
 g(t,z)=
 \begin{cases}
  \sigma(t)q_t^2G(z/q_t)+G(-z),&q_t>0,\\
  G(-z),&q_t=0.
 \end{cases}
\end{equation}
The coefficient $q$ was fixed before choosing a solution parameter. Hence this is one predictable generator for all the solutions below. It does not depend on $y$.

The two terms in~\eqref{generator} have disjoint positive and negative supports, so $|g(t,z)|\leq z^2$ and $g(t,0)=0$. At points of differentiability, the scaled tent has derivative of magnitude at most $2q_t$, supported where $z\geq q_t/2$; its derivative is therefore bounded by $4|z|$. The other tent has derivative bounded by $2$. Integration along the interval from $z$ to $z'$ proves~\eqref{bounds}. Each $g(t,\omega,\cdot)$ is compactly supported and nonzero on the negative axis. A nonzero compactly supported function on $\mathbb R$ can be neither convex nor concave.

For $c\in[0,1]$ put
\[
 Y^c=X+cD,\qquad Z^c=q+cK,\qquad \xi=X_T.
\]
Where $q>0$, equations~\eqref{D} give
\[
 \frac{Z^c}{q}\in[13/15,17/15]\subset[1/2,3/2],
 \qquad
 g(t,Z^c)=\sigma q^2+2c\sigma qK.
\]
Where $q=0$, both $Z^c$ and $g(t,Z^c)$ vanish. Combining~\eqref{q} and~\eqref{linear} now yields
\[
 dY^c_t=Z^c_t\,dB_t-g(t,Z^c_t)\,dt.
\]
Since $D_T=0$, integration gives~\eqref{bsde} with the same $\xi$ for every $c$. The identities hold simultaneously for all times by continuity, including $T$, since only finitely many legs are used almost surely.

Finally, $|Z^c|\leq(17/15)q$ and $\sup|Y^c|\leq2a_N+\eta$. Equations~\eqref{tails}--\eqref{maximum} prove~\eqref{moments}. Each stopped exponential is dominated by the integrable random variable $e^{\mu\sup|Y^c|}$, which proves class (D). The initial values $Y^c_0=c\eta d_1$ are distinct. This proves Theorem~\ref{main}.\hfill$\square$

\paragraph{Scope.}
The generator is random and time-dependent, as permitted in the stated problem. Taking $g(t,y,z)=g(t,z)$ satisfies any positive uniform Lipschitz bound in $y$. The example therefore meets the repaired hypotheses with $\delta=1$ and $C=4$; it does not use the missing-$y$-regularity obstruction. In higher Brownian dimension, use the first coordinate and replace $g(t,z)$ by $g(t,z_1)$.

\begin{thebibliography}{9}
\bibitem{FHT} S. Fan, Y. Hu and S. Tang.
\emph{1D nonlinear backward stochastic differential equations: a unified theory and applications}.
arXiv:2309.06233v2, 11 February 2026, \S1, equation (3.12), and Problem~5.5.
\url{https://arxiv.org/html/2309.06233v2}.
\end{thebibliography}
\smallskip\noindent\footnotesize
The construction and manuscript were prepared with Codex assistance. The accompanying records specify the source comparison and mathematical checks.

\normalsize
\section*{Submission scope and verification}
This is a complete counterexample to the repaired C65-29 question. The earlier catalogue counterexample and Pro solved badge concern the literal statement's missing condition in $y$; both prior manuscripts explicitly leave the general repaired question unresolved. The present generator is independent of $y$, has zero value at the origin, satisfies the uniform quadratic Lipschitz bound, and admits distinct solutions in every required exponential class (D). Thus this submission resolves the remaining repaired question and also refutes the literal statement.

Author: Ji Ho Bae. Prepared and analytically checked through GPT Codex; the exact serving revision was not exposed. Two full analytic passes were performed by the authoring assistant. No new independent-agent review, external peer review, or Lean verification is claimed.

\begin{itemize}
\item \href{https://github.com/jbaelaw/nonconvex-bsde-all-exponential-nonuniqueness/blob/841d20e8862a2ab7bac7d7b212088c21056e8fd6/paper/Proof.pdf}{Complete three-page proof PDF} and \href{https://github.com/jbaelaw/nonconvex-bsde-all-exponential-nonuniqueness/blob/841d20e8862a2ab7bac7d7b212088c21056e8fd6/paper/Proof.tex}{TeX source}.
\item \href{https://github.com/jbaelaw/nonconvex-bsde-all-exponential-nonuniqueness/blob/841d20e8862a2ab7bac7d7b212088c21056e8fd6/verification/Review.md}{Mathematical checks} and \href{https://github.com/jbaelaw/nonconvex-bsde-all-exponential-nonuniqueness/blob/841d20e8862a2ab7bac7d7b212088c21056e8fd6/verification/SourceAudit.md}{source and prior-work scope}.
\item \href{https://github.com/jbaelaw/nonconvex-bsde-all-exponential-nonuniqueness/blob/841d20e8862a2ab7bac7d7b212088c21056e8fd6/verification/Verification.json}{Verification record} and \href{https://github.com/jbaelaw/nonconvex-bsde-all-exponential-nonuniqueness/blob/841d20e8862a2ab7bac7d7b212088c21056e8fd6/provenance/Process.md}{generation process}.
\end{itemize}
The canonical statement, original rank and prior attempts are retained. Submitted for expert review; no catalogue-maintainer acceptance is claimed.

\end{document}
